% =====================================================================
%  ThmSmith Stage -1 proposal skeleton.
%  Written automatically by the ThmSmith TS pipeline before Stage -1 runs.
%  Locked parameters: qid=eid\_lp\_svar\_nonequiv, specialization=v1
%
%  HOW TO USE THIS FILE
%  --------------------
%  - The preamble (everything above the `% --- BEGIN CONTENT ---` marker)
%    and the `\section{...}` headers below are fixed by the pipeline.
%    Do NOT rewrite or rename them; the Step 3 reviewer subagent and the
%    downstream Stage 0 / Stage 1 prompts grep for these section titles.
%  - Each section contains exactly one `% TODO(stage-1 ...): ...`
%    placeholder. Replace each TODO with the content described for that
%    section in `ThmSmith/tools/src/prompts/stage_neg1_draft.txt`
%    ("REQUIRED OUTPUT FORMAT (Stage -1 proposal .tex)").
%  - The `\paragraph{Locked parameters.}` block in the metadata block
%    must pin the chosen `(qid, specialization, p, assignment, target,
%    regression)` precisely. If the proposal maps to the Q1 pool-mode,
%    reproduce that block verbatim from
%    `ThmSmith/doc/panel_formula_question_generator_note_with_common_prompts.tex`
%    (Q2..Q6 templates have been retired). Otherwise (free-form
%    `(qid, specialization)` — the default), define each parameter
%    explicitly here (no implicit conventions). Stage 0 quotes this block;
%    do not rephrase it after locking.
%  - Removing a section header, renaming a macro, or rewriting the
%    preamble is a regression: the file must still match this skeleton
%    after you finish.
% =====================================================================

\documentclass[11pt]{article}

\usepackage{amsmath,amssymb,amsthm,mathtools}
\usepackage{enumitem}
\usepackage[colorlinks=true,linkcolor=blue,citecolor=blue,urlcolor=blue]{hyperref}
\usepackage{natbib}

\newcommand{\E}{\mathbb{E}}
\renewcommand{\Pr}{\mathbb{P}}
\newcommand{\one}{\mathbf{1}}
\newcommand{\transpose}{\mathsf{T}}
\newcommand{\calH}{\mathcal{H}}
\newcommand{\calF}{\mathcal{F}}
\newcommand{\calR}{\mathcal{R}}
\newcommand{\R}{\mathbb{R}}
\newcommand{\plim}{\operatorname*{plim}}

\theoremstyle{definition}
\newtheorem{definition}{Definition}
\newtheorem{assumption}{Assumption}
\newtheorem{example}{Example}

\theoremstyle{plain}
\newtheorem{theorem}{Theorem}
\newtheorem{conjecture}{Conjecture}
\newtheorem{proposition}{Proposition}
\newtheorem{lemma}{Lemma}
\newtheorem{corollary}{Corollary}

\theoremstyle{remark}
\newtheorem{remark}{Remark}

\title{ThmSmith Stage -1 proposal: \texttt{eid\_lp\_svar\_nonequiv} / \texttt{v1}%
  \\ \large\textit{A common-cumulant boundary for LP/SVAR impulse-response non-equivalence}}
\author{ThmSmith Stage -1 proposer}
\date{\today}

\begin{document}
\maketitle

% --- BEGIN CONTENT ---  (everything above is fixed by the pipeline)

\paragraph{Locked parameters.}
This is a partial-identification / linear-projection anchor, not a Q1 pool-mode.  The locked tuple is
\[
\begin{gathered}
(\texttt{qid},\texttt{specialization},P,\theta,\text{identifying restrictions},
  \text{bounding functional})\\
=
(\texttt{eid\_lp\_svar\_nonequiv},\texttt{v1},
  P_{\{Y_t\}_{t\in\mathbb Z}}^{(4,H)},\
  \theta_{j,0:H}^{\mathrm{SIRF}},\
  \mathcal R_{\mathrm{VAR},H}^{\mathrm{NG+sign}},\
  [\Theta_{\mathrm{SVAR}}(P),\Theta_{\mathrm{LP}}^{\mathrm{sep}}(P),
   \Delta_{\mathrm{NG}}(P)]) .
\end{gathered}
\]
Here \(P_{\{Y_t\}_{t\in\mathbb Z}}^{(4,H)}\) is the stationary observable law of the finite block needed to recover the reduced-form innovation covariance \(\Sigma_u\), the LP population rows \(M_h=\E[Y_{t+h}u_t^\transpose]\Sigma_u^{-1}\) for \(0\le h\le H\), and the whitened fourth-cumulant tensor \(T_4=\operatorname{cum}_4(\Sigma_u^{-1/2}u_t)\).  The parameter \(\theta_{j,0:H}^{\mathrm{SIRF}}\) is the horizon-\(0\) through horizon-\(H\) structural impulse-response path to a labeled structural shock \(j\).  The restrictions \(\mathcal R_{\mathrm{VAR},H}^{\mathrm{NG+sign}}\) are: a finite-order linear reduced-form VAR representation for the LP/SVAR population rows; an impact representation \(u_t=A\varepsilon_t\) with \(AA^\transpose=\Sigma_u\), mutually independent unit-variance shocks, finite fourth moments, and fourth cumulants \(\kappa_j\) satisfying a stated non-Gaussian gap bound; partial zero/sign restrictions on the impact vector and selected response horizons; and a fixed labeling convention for the shock of interest.  The bounding functional records the sharp common-law SVAR identified set \(\Theta_{\mathrm{SVAR}}(P)\), the horizonwise LP relaxation \(\Theta_{\mathrm{LP}}^{\mathrm{sep}}(P)\), and the non-equivalence gap functional \(\Delta_{\mathrm{NG}}(P)=\sup_{\theta\in\Theta_{\mathrm{LP}}^{\mathrm{sep}}(P)}\operatorname{dist}(\theta,\Theta_{\mathrm{SVAR}}(P))\).

\begin{abstract}
Plagborg-Moller and Wolf show that local projections and VARs estimate the same population impulse responses when they target the same reduced-form object and use the same structural restrictions.  Non-Gaussian SVAR identification changes the object: the admissible impact matrix must diagonalize one shared innovation-law cumulant tensor, whereas horizonwise LP restrictions can be read as separate projection restrictions unless a common-shock completion is imposed.  This proposal asks for the sharp algebraic boundary between those two identified sets under bounded fourth-cumulant gaps and partial sign/zero restrictions.  The routine theorem is the Gaussian and common-cumulant-matching reduction back to the Plagborg-Moller--Wolf equivalence set.  The flagship conjecture is that strict LP/SVAR non-equivalence holds on an open set exactly when the horizonwise LP rotations cannot be pasted into one orthogonal rotation that diagonalizes the same fourth-cumulant tensor.  Resolving it would tell applied macroeconomists when a non-Gaussian or independence-refined sign-restricted SVAR is not just an LP implementation choice, but a strictly sharper structural object.
\end{abstract}

\section{Introduction}
\paragraph{Why the problem is important.}
Applied structural macro work now routinely moves between LPs, VARs, proxy SVARs, sign restrictions, and non-Gaussian statistical identification.  The LP/VAR equivalence theorem says this move is harmless when the identifying restrictions are attached to the same population impulse-response object.  Independence-based non-Gaussian SVARs are different because the restriction is not merely a row restriction on a response path: it is a one-law restriction on the entire reduced-form innovation vector.

\paragraph{The main finding (proposed).}
The proposed headline result is Conjecture~\ref{conj:common-cumulant-boundary}: the LP relaxation and the non-Gaussian sign-restricted SVAR identified set coincide if and only if the LP-admissible horizonwise rotations satisfy a common-cumulant matching condition.  If that condition fails with a fourth-cumulant gap at least \(\delta_\kappa>0\), then the LP set strictly contains response paths that no independent non-Gaussian SVAR shock law can generate.

\paragraph{What is new.}
The closest published result is \citet{PlagborgMollerWolf2021}, whose epsilon-gap is that its portability claim holds for common reduced-form impulse-response estimands and standard VAR identification schemes, while it does not characterize when higher-cumulant independence restrictions force one shared structural completion across all horizons.  The closest non-Gaussian SVAR papers, \citet{GourierouxMonfortRenne2017}, \citet{LanneMeitzSaikkonen2017}, and \citet{DrautzburgWright2023}, identify or refine SVAR rotations using non-Gaussianity and independence, but they do not compare the resulting common-law set with a horizonwise LP projection relaxation.

\paragraph{How it positions in the literature.}
This is a panel / linear-projection methods proposal at the LP/SVAR frontier, with a partial-identification identified-set kernel.  It uses the LP/SVAR equivalence literature as the baseline, sign-restricted SVARs as the partial-ID geometry, and non-Gaussian SVAR identification as the common-shock completion restriction.  It is not an SCM or do-calculus proposal.

\paragraph{Implications for the community.}
The direct consumer is \citet{DrautzburgWright2023}: their independence-refined sign-restricted VAR sets could be reported with a sharp diagnostic saying whether the same refinement is reproducible by horizonwise LP restrictions or whether a common non-Gaussian shock law strictly tightens the response set.

\paragraph{How research benefits from the conclusion.}
Researchers should report a common-cumulant compatibility check before treating non-Gaussian/sign-refined SVAR conclusions as interchangeable with LP-based sign-restricted impulse-response analysis.

\section{Related work}
\citet{Jorda2005} introduced LP impulse-response estimation and supplies the projection side of the comparison.  \citet{PlagborgMollerWolf2021} prove that LPs and VARs estimate the same population impulse responses under unrestricted lag structures and that standard VAR identification schemes, including sign restrictions, can be ported to LPs.  \citet{MontielOleaPlagborgMoller2021} and \citet{MontielOleaPlagborgMollerQianWolf2024} show why LP robustness to persistence and lag misspecification is now a live methodological advantage rather than a textbook alternative.  \citet{LiPlagborgMollerWolf2024} document the practical bias-variance frontier between LP and VAR estimators across many DGPs, which motivates separating finite-sample estimator differences from population identified-set differences.

\citet{GourierouxMonfortRenne2017} and \citet{LanneMeitzSaikkonen2017} are the core non-Gaussian SVAR anchors: they use independent non-Gaussian shocks to remove the usual rotation non-identification up to scale, permutation, and labeling restrictions.  \citet{Guay2021} gives GMM local identification conditions for non-Gaussian SVARs, making rank and moment conditions explicit.  \citet{DrautzburgWright2023} is the closest identified-set paper because it refines sign-restricted VAR sets using independence and notes non-convexity and weak higher-moment identification.  \citet{LanneLiu2026} shows that orthogonality plus higher moments can identify some shocks without mutual independence, marking the current frontier around exactly which common-law restrictions are required.

\citet{BaumeisterHamilton2015} and \citet{AriasRubioRamirezWaggoner2018} supply the sign-restricted SVAR background: the former explains prior-driven underidentification, and the latter gives algorithms and theory for sign/zero restricted SVAR inference.  \citet{MontielOleaStockWatson2021} is relevant because weak external-instrument SVAR inference already treats impulse responses as set-valued under weak identification, a cousin of the near-Gaussian weak-cumulant regime considered here.

The in-catalogue prior art is adjacent but not duplicative.  The ThmSmith API records Q1 minimal-basis and sign-flip panel theorems about FWL contamination under binary treatment histories, but no LP/SVAR or non-Gaussian innovation theorem.  The closest banked proposal family \texttt{q1\_spectral\_phase\_transition} studies thresholded row-space identification under Markov assignment; its headline kernel is a panel event-window singular-subspace boundary, not a structural time-series cumulant-completion boundary.  The closest partial-ID templates \texttt{pid\_network\_exposure}, \texttt{pid\_dose\_response\_iv}, and \texttt{pid\_offline\_policy\_joint\_misspec} study shared-completion or non-rectangularity gaps in other causal settings; they motivate the common-completion form but do not mention Plagborg-Moller--Wolf, SVARs, structural impulse responses, or non-Gaussian innovations.

\section{Formal setup}
Let \(Y_t\in\R^n\) be a stationary mean-zero time series with reduced-form innovation \(u_t\), innovation covariance \(\Sigma_u\succ0\), and population LP rows
\[
M_h=\E[Y_{t+h}u_t^\transpose]\Sigma_u^{-1},\qquad h=0,\ldots,H,
\]
where \(M_0=I_n\).  Write \(v_t=\Sigma_u^{-1/2}u_t\) and let \(T_4=\operatorname{cum}_4(v_t)\) be the symmetric fourth-cumulant tensor.  A structural completion is an orthogonal matrix \(Q\in\mathcal O(n)\), an impact matrix \(A=\Sigma_u^{1/2}Q\), and shocks \(\varepsilon_t=Q^\transpose v_t\).  The response path for labeled shock \(j\) is
\[
\theta_j(Q)=\bigl(M_h\Sigma_u^{1/2}Qe_j\bigr)_{h=0}^H.
\]

The common-law SVAR identified set is
\[
\Theta_{\mathrm{SVAR}}(P)
=
\{\theta_j(Q): Q\in\mathcal O(n),\ Q\text{ satisfies the sign/zero restrictions, }
T_4(Q)\text{ is diagonal with gap }\delta_\kappa\},
\]
where \(T_4(Q)=T_4\times_1 Q^\transpose\times_2 Q^\transpose\times_3 Q^\transpose\times_4 Q^\transpose\).  The horizonwise LP relaxation is
\[
\Theta_{\mathrm{LP}}^{\mathrm{sep}}(P)
=
\left\{\bigl(M_h\Sigma_u^{1/2}q_h\bigr)_{h=0}^H:
q_h\in\mathbb S^{n-1}\text{ satisfies the horizon }h\text{ sign/zero and marginal cumulant bounds}\right\},
\]
where \(q_h\) may vary with \(h\).  The common-cumulant matching restriction is the event that every LP-admissible path has a representation \(q_h=Qe_j\) for one orthogonal \(Q\) that diagonalizes \(T_4\) and uses one shock label \(j\) across all horizons.

\begin{quote}\small
\begin{verbatim}
(qid, specialization) =
  (eid_lp_svar_nonequiv, v1)
anchor cluster =
  partial-identification / panel-linear-projection frontier
observable distribution P =
  stationary finite-fourth-moment law of {Y_t} sufficient to recover
  Sigma_u, M_h for h=0,...,H, and T_4=cum_4(Sigma_u^{-1/2}u_t)
parameter theta =
  theta_{j,0:H}^{SIRF}=(M_h Sigma_u^{1/2} Q e_j)_{h=0}^H
identifying restrictions =
  finite-order linear reduced-form VAR rows M_h; structural completion
  u_t=A eps_t, AA'=Sigma_u; independent unit-variance shocks with fourth
  cumulants kappa_j and gap delta_kappa; partial zero/sign restrictions;
  fixed shock-label convention; and common-cumulant matching when imposed
bounding functional =
  Theta_SVAR(P), Theta_LP^sep(P), and
  Delta_NG(P)=sup_{theta in Theta_LP^sep(P)}
    dist(theta,Theta_SVAR(P))
\end{verbatim}
\end{quote}

\section{Assumptions}
\begin{assumption}[Reduced-form LP/VAR population rows]\label{ass:rf-rows}
The observable law \(P\) is stationary, has finite fourth moments, has \(\Sigma_u\succ0\), and determines \(M_h=\E[Y_{t+h}u_t^\transpose]\Sigma_u^{-1}\) for every \(h=0,\ldots,H\).
\end{assumption}

\begin{assumption}[Independent non-Gaussian structural completion]\label{ass:ng-completion}
An admissible SVAR completion has \(u_t=A\varepsilon_t\), \(AA^\transpose=\Sigma_u\), \(\E[\varepsilon_t]=0\), \(\E[\varepsilon_t\varepsilon_t^\transpose]=I_n\), mutually independent shock components, finite fourth moments, and fourth cumulants \(\kappa_1,\ldots,\kappa_n\) with non-Gaussian gap \(\delta_\kappa=\min_{r\ne s}|\kappa_r-\kappa_s|\) on the labeled components used for identification.
\end{assumption}

\begin{assumption}[Partial sign/zero label restrictions]\label{ass:sign-zero}
For a fixed labeled shock \(j\), the admissible response path must satisfy a finite list of homogeneous zero restrictions and weak sign restrictions on entries of \(M_h\Sigma_u^{1/2}Qe_j\) for \(h\in\mathcal H_S\subseteq\{0,\ldots,H\}\), including at least one nonzero label-normalizing inequality.
\end{assumption}

\begin{assumption}[LP separated relaxation]\label{ass:lp-sep}
The LP comparator \(\Theta_{\mathrm{LP}}^{\mathrm{sep}}(P)\) imposes the restrictions in Assumption~\ref{ass:sign-zero} horizon by horizon and imposes only the marginal fourth-cumulant bounds \(\kappa_j\in[\underline\kappa_j,\overline\kappa_j]\); it does not require the vectors \(q_h\) to be columns of one orthogonal matrix that diagonalizes \(T_4\).
\end{assumption}

\begin{assumption}[Common-cumulant matching]\label{ass:cumulant-match}
The common-cumulant matching condition holds for a response path if there exist one \(Q\in\mathcal O(n)\), one label \(j\), and one cumulant vector \(\kappa\) such that \(q_h=Qe_j\) for all horizons represented in the path and \(T_4(Q)_{rstu}=0\) unless \(r=s=t=u\), with diagonal entries \(T_4(Q)_{rrrr}=\kappa_r\).
\end{assumption}

\begin{assumption}[Nondegenerate witness region]\label{ass:witness-region}
The reduced-form rows \((M_h)_{h=0}^H\), the sign/zero cone, and the tensor \(T_4\) lie in an open region where the labeled response path is not annihilated by all \(M_h\Sigma_u^{1/2}\), the sign inequalities used for labeling are locally stable, and at least one LP-admissible horizon vector violates Assumption~\ref{ass:cumulant-match}.
\end{assumption}

\section{Main results}
\begin{theorem}[Theorem 1: Gaussian and matched-cumulant reduction]\label{thm:gaussian-reduction}
Under Assumptions~\ref{ass:rf-rows}, \ref{ass:sign-zero}, and \ref{ass:lp-sep}, if either \(T_4=0\) or Assumption~\ref{ass:cumulant-match} is imposed as part of the LP restrictions, then
\[
\Theta_{\mathrm{LP}}^{\mathrm{sep}}(P)=\Theta_{\mathrm{SVAR}}(P)
\]
after replacing the non-Gaussian diagonalization restriction by the same sign/zero restrictions used in both procedures.  In this case \(\Delta_{\mathrm{NG}}(P)=0\), and the common population response set is the Plagborg-Moller--Wolf LP/VAR identified set for the stated restrictions.
\end{theorem}
\paragraph{Why this is new and worth studying.}
(a) This is a calibration theorem, not the novelty kernel.  Its closest prior art is \citet{PlagborgMollerWolf2021}; the gap is that the statement pins the exact corner in which the present proposal must reduce to their equivalence result.  (b) The concrete consumer is \citet{LiPlagborgMollerWolf2024}: the theorem separates their estimator bias-variance comparison from any population non-equivalence caused by non-Gaussian structural restrictions.

\begin{conjecture}[Conjecture 1: common-cumulant boundary for LP/SVAR non-equivalence (NOT YET PROVED)]\label{conj:common-cumulant-boundary}
Under Assumptions~\ref{ass:rf-rows}--\ref{ass:witness-region}, the following are equivalent:
\[
\Theta_{\mathrm{LP}}^{\mathrm{sep}}(P)=\Theta_{\mathrm{SVAR}}(P)
\quad\Longleftrightarrow\quad
\text{every extreme LP-admissible response path satisfies Assumption~\ref{ass:cumulant-match}.}
\]
If the right-hand condition fails and \(\delta_\kappa>0\), then the containment is strict on an open set:
\[
\Theta_{\mathrm{SVAR}}(P)\subsetneq\Theta_{\mathrm{LP}}^{\mathrm{sep}}(P),
\qquad
\Delta_{\mathrm{NG}}(P)>0.
\]
Moreover, for fixed \((\Sigma_u,M_0,\ldots,M_H)\) and locally stable sign/zero labels, the gap admits a closed-form algebraic certificate: a nonzero mixed entry of \(T_4(Q)\) or a violated diagonal-cumulant equality for every orthogonal \(Q\) whose column realizes the LP horizon vectors.
\end{conjecture}
\paragraph{Why this is new and worth studying.}
(a) The closest published results are \citet{PlagborgMollerWolf2021}, \citet{GourierouxMonfortRenne2017}, \citet{LanneMeitzSaikkonen2017}, and \citet{DrautzburgWright2023}.  The epsilon-gap is that none states the necessary-and-sufficient pasting condition under which horizonwise LP sign/cumulant restrictions are generated by one independent non-Gaussian SVAR innovation law.  (b) The concrete consumer is the independence-refined sign-restricted workflow in \citet{DrautzburgWright2023}: the result would provide a reportable certificate for when their non-Gaussian refinement is genuinely sharper than an LP horizon-by-horizon relaxation.

\begin{conjecture}[Conjecture 2: two-shock rational witness and sharpness certificate (NOT YET PROVED)]\label{conj:two-shock-witness}
Under Assumptions~\ref{ass:rf-rows}--\ref{ass:witness-region}, there exists a two-variable, two-horizon rational observable law \(P^\star\) with positive definite \(\Sigma_u^\star\), bounded fourth cumulants, \(\delta_\kappa^\star>0\), and nondegenerate partial sign/zero restrictions such that
\[
\Theta_{\mathrm{SVAR}}(P^\star)\subsetneq\Theta_{\mathrm{LP}}^{\mathrm{sep}}(P^\star).
\]
The witness can be chosen so that both sets have nonempty relative interior in their natural response coordinates, and strictness is certified by a single polynomial inequality in the whitened fourth-cumulant tensor entries.  In the limit \(\delta_\kappa^\star\downarrow0\), the polynomial certificate vanishes and the witness collapses to the Gaussian/sign-only equivalence corner of Theorem~\ref{thm:gaussian-reduction}.
\end{conjecture}
\paragraph{Why this is new and worth studying.}
(a) \citet{BaumeisterHamilton2015} and \citet{AriasRubioRamirezWaggoner2018} explain sign-restricted SVAR set geometry, while \citet{LanneMeitzSaikkonen2017} and \citet{DrautzburgWright2023} explain non-Gaussian refinement; the epsilon-gap is an explicit nondegenerate observable witness separating the LP relaxation from the common-law SVAR set.  (b) The concrete consumer is \citet{MontielOleaPlagborgMollerQianWolf2024}: a small rational witness would show that lag-misspecification-robust LP inference and non-Gaussian SVAR identification can target different population sets before any sampling issue appears.

\section{Sanity-check corollaries}
\begin{corollary}[Gaussian sign-only collapse]\label{cor:gaussian}
Under Assumptions~\ref{ass:rf-rows}, \ref{ass:sign-zero}, and \ref{ass:lp-sep}, if \(T_4=0\) and no non-Gaussian cumulant gap is imposed, Theorem~\ref{thm:gaussian-reduction} collapses to the sign/zero LP/VAR equivalence result of \citet{PlagborgMollerWolf2021}: both procedures optimize the same response vectors \(M_h\Sigma_u^{1/2}q\) over the same sign/zero cone.
\end{corollary}

\begin{corollary}[No artificial gap under common matching]\label{cor:matched}
Under Assumptions~\ref{ass:rf-rows}--\ref{ass:cumulant-match}, every response path in \(\Theta_{\mathrm{LP}}^{\mathrm{sep}}(P)\) is generated by one diagonalizing \(Q\), so the proposed gap functional is zero by set equality, not by a normalization convention.
\end{corollary}

\begin{corollary}[Witness nondegeneracy check]\label{cor:witness-nondegenerate}
The rational witness required by Conjecture~\ref{conj:two-shock-witness} must keep \(\Sigma_u^\star\succ0\), positive sign-restriction slack for at least one label inequality, and \(\delta_\kappa^\star>0\).  Therefore strictness cannot be witnessed by a collapsed innovation law, a zero-response horizon, or a hand-tuned sign boundary.
\end{corollary}

\section{Proof-sketch route}
The mathematical route is finite-dimensional.  Whiten the innovations, express every SVAR completion as \(A=\Sigma_u^{1/2}Q\), and write non-Gaussian independence as orthogonal diagonalization of the symmetric fourth-cumulant tensor \(T_4\).  The LP relaxation is a product or projection of horizonwise sign/zero cones through the linear maps \(q\mapsto M_h\Sigma_u^{1/2}q\).  The equality boundary should follow by comparing extreme points of the LP relaxation with the algebraic variety cut out by the mixed-cumulant equations \(T_4(Q)_{rstu}=0\) for non-diagonal indices; strictness is certified by a violated polynomial plus sign slack.  The small witness can be found by enumerating a two-shock orthogonal rotation angle, two response rows, and a diagonal fourth-cumulant vector, then checking positivity and sign slack by explicit polynomial inequalities.

\section{Honest scope flags}
Theorem~\ref{thm:gaussian-reduction} is intended as a routine calibration theorem.  Conjectures~\ref{conj:common-cumulant-boundary} and \ref{conj:two-shock-witness} are not yet proved; the exact form of the closed-form certificate may need to be a support-function certificate rather than a single polynomial in higher dimensions.  The proposal deliberately uses a horizonwise LP relaxation, not the fully restricted LP implementation that already imposes the same common structural law as the SVAR; without that distinction the conjecture would contradict \citet{PlagborgMollerWolf2021}.  No claim is made about finite-sample coverage, Bayesian priors over rotations, or noncausal VAR models.

\section{Iteration trail}
\begin{itemize}[leftmargin=*]
\item v1 (cold-start) --- initial draft from Stage -1.1 common-shock-completion gap; prior verdict: none.
\end{itemize}

\section*{Bibliography}
\begin{thebibliography}{99}
\bibitem[Arias, Rubio-Ramirez, and Waggoner(2018)]{AriasRubioRamirezWaggoner2018}
Arias, J. E., J. F. Rubio-Ramirez, and D. F. Waggoner (2018).
Inference based on structural vector autoregressions identified with sign and zero restrictions: Theory and applications.
\emph{Econometrica} 86(2), 685--720.

\bibitem[Baumeister and Hamilton(2015)]{BaumeisterHamilton2015}
Baumeister, C. and J. D. Hamilton (2015).
Sign restrictions, structural vector autoregressions, and useful prior information.
\emph{Econometrica} 83(5), 1963--1999.

\bibitem[Drautzburg and Wright(2023)]{DrautzburgWright2023}
Drautzburg, T. and J. H. Wright (2023).
Refining set-identification in VARs through independence.
\emph{Journal of Econometrics} 235(2), 1827--1847.

\bibitem[Gourieroux, Monfort, and Renne(2017)]{GourierouxMonfortRenne2017}
Gourieroux, C., A. Monfort, and J.-P. Renne (2017).
Statistical inference for independent component analysis: Application to structural VAR models.
\emph{Journal of Econometrics} 196(1), 111--126.

\bibitem[Guay(2021)]{Guay2021}
Guay, A. (2021).
GMM estimation of non-Gaussian structural vector autoregression.
\emph{Journal of Business \& Economic Statistics} 39(1), 34--49.

\bibitem[Jorda(2005)]{Jorda2005}
Jorda, O. (2005).
Estimation and inference of impulse responses by local projections.
\emph{American Economic Review} 95(1), 161--182.

\bibitem[Lanne and Liu(2026)]{LanneLiu2026}
Lanne, M. and K. Liu (2026).
Identifying structural vector autoregressions via non-Gaussianity of potentially dependent shocks.
\emph{The Econometrics Journal}, corrected proof.

\bibitem[Lanne, Meitz, and Saikkonen(2017)]{LanneMeitzSaikkonen2017}
Lanne, M., M. Meitz, and P. Saikkonen (2017).
Identification and estimation of non-Gaussian structural vector autoregressions.
\emph{Journal of Econometrics} 196(2), 288--304.

\bibitem[Li, Plagborg-Moller, and Wolf(2024)]{LiPlagborgMollerWolf2024}
Li, D., M. Plagborg-Moller, and C. K. Wolf (2024).
Local projections vs. VARs: Lessons from thousands of DGPs.
\emph{Journal of Econometrics} 244, 105722.

\bibitem[Montiel Olea and Plagborg-Moller(2021)]{MontielOleaPlagborgMoller2021}
Montiel Olea, J. L. and M. Plagborg-Moller (2021).
Local projection inference is simpler and more robust than you think.
\emph{Econometrica} 89(4), 1789--1823.

\bibitem[Montiel Olea, Plagborg-Moller, Qian, and Wolf(2024)]{MontielOleaPlagborgMollerQianWolf2024}
Montiel Olea, J. L., M. Plagborg-Moller, E. Qian, and C. K. Wolf (2024).
Double robustness of local projections and some unpleasant VARithmetic.
NBER Working Paper 32495; \emph{Econometrica}, forthcoming.

\bibitem[Montiel Olea, Stock, and Watson(2021)]{MontielOleaStockWatson2021}
Montiel Olea, J. L., J. H. Stock, and M. W. Watson (2021).
Inference in structural vector autoregressions identified with an external instrument.
\emph{Journal of Econometrics} 225(1), 74--87.

\bibitem[Plagborg-Moller and Wolf(2021)]{PlagborgMollerWolf2021}
Plagborg-Moller, M. and C. K. Wolf (2021).
Local projections and VARs estimate the same impulse responses.
\emph{Econometrica} 89(2), 955--980.
\end{thebibliography}

\end{document}
